% !TeX root = ../main.tex

% ==================================================
% 2.2. 클래스도 객체인가? - class 정의를 실행하면 무엇이 생성되는가?
% ==================================================

\subsection{클래스도 객체인가? - class 정의를 실행하면 무엇이 생성되는가?}

이번에는 Python의 \zsh{class}문을 단순히 클래스를 선언하는 문법으로 이해하는 것이 아니라, Python이 \zsh{class}문을 실행할 때 어떤 객체를 생성하고, 그 객체를 어떤 이름에 바인딩하는지 살펴본다.

함수와 데코레이터에서 다음 개념을 학습했을 것이다.

\begin{dhkcode}
def add(a, b):
    return a + b
\end{dhkcode}

이 코드를 실행하면 Python은 함수 객체를 생성하고 \zsh{add}라는 이름에 바인딩한다.

클래스도 비슷한 관점에서 이해할 수 있다.

\begin{dhkcode}
class Person:
    pass
\end{dhkcode}

이 코드를 실행하면 Python은 \zsh{클래스 객체(Class Object)}를 생성하고 \zsh{Person}이라는 이름에 바인딩한다.

다만 함수 객체와 클래스 객체는 서로 다른 종류의 객체이며, 생성 과정과 역할에도 차이가 있다.

이번에는 그 차이를 순서대로 살펴본다.

\subsubsection{개념 정의 - Python에서 클래스도 객체다}

Python에서는 정수, 문자열, 리스트, 함수뿐 아니라 클래스도 객체로 취급한다.

다음 코드를 살펴보자.

\begin{dhkcode}
class Person:
    pass

print(Person)
print(type(Person))
\end{dhkcode}

실행 결과:

\begin{dhkzsh}
<class '__main__.Person'>
<class 'type'>
\end{dhkzsh}

첫 번째 출력은 \zsh{Person}이라는 이름이 클래스 객체를 참조하고 있다는 사실을 보여준다.

두 번째 출력은 해당 클래스 객체의 타입이 \zsh{type}이라는 사실을 보여준다.

여기서 중요한 구분이 있다.

\begin{dhknote}
\begin{dhksmry}[centered]
클래스 객체(class object)
\end{dhksmry}
\end{dhknote}

\begin{dhkcode}
Person
\end{dhkcode}

\zsh{Person}은 클래스 객체를 참조하는 이름이다.

\begin{dhknote}
\begin{dhksmry}[centered]
인스턴스 객체(Instance Object)
\end{dhksmry}
\end{dhknote}

\begin{dhkcode}
Person()
\end{dhkcode}

\zsh{Person()}은 클래스 객체를 호출하는 표현식이다.

기본적인 클래스에서는 이 호출을 통해 새로운 인스턴스가 생성된다.

\begin{dhknote}
\begin{dhksmry}[centered]
인스턴스를 참조하는 이름
\end{dhksmry}
\end{dhknote}

\begin{dhkcode}
p = Person()
\end{dhkcode}

생성된 인스턴스 객체를 \zsh{p}라는 이름에 바인딩한다.

따라서 다음과 같이 구분할 수 있다.

\begin{dhkdiagram}[in=10]
node p (1,1) :: Person
node c (3,1) :: 클래스 객체
node i (5,1) :: 인스턴스 객체
node s (7,1) :: p

edge p -> c
edge c -> i :: Person() 호출
edge s -> i
\end{dhkdiagram}

여기서 \zsh{Person}과 \zsh{p}는 모두 이름이지만, 참조하는 객체의 종류가 다르다.

\begin{dhktable}
|표현식|	의미|	참조하거나 반환하는 객체|
|:---|:---|:---|
|Person|	클래스 이름을 평가|	클래스 객체|
|Person()|	클래스 객체 호출|	일반적으로 새로운 //*인스턴스 객체|
|p = Person()|	반환된 객체를//*이름에 바인딩|	p가 인스턴스를 참조|
|type(Person)|	클래스 객체의 //*타입 확인|	type|
|type(p)|	인스턴스 객체의 //*타입 확인|	Person|
\end{dhktable}

이 관계를 실제 코드로 확인해 보자.

\begin{dhkcode}
class Person:
    pass

p = Person()

print(type(Person))
print(type(p))

print(isinstance(p, Person))
\end{dhkcode}

실행 결과:

\begin{dhkzsh}
<class 'type'>
<class '__main__.Person'>
True
\end{dhkzsh}

여기서 \zsh{isinstance(p, Person)}은 \zsh{p}가 \zsh{Person} 클래스의 인스턴스인지 확인하는 표현식이다.

상속 관계까지 고려하므로 단순히 \zsh{type(p) is Person}과 항상 같은 검사는 아니다.

핵심은 다음과 같다.

\begin{dhksmry}[centered]
Python에서 클래스는 객체를 만들기 위한 정의이면서 동시에, 

Python 프로그램 안에서 참조하고 전달할 수 있는 객체이기도 하다.
\end{dhksmry}

\subsubsection{\zsh{class}문은 단순한 선언이 아니라 실행되는 문장이다}

이 부분이 이번 학습에서 가장 중요하다.

Python의 \zsh{class}문은 프로그램 실행 과정에서 평가되고 실행된다.

예를 들어 다음 코드를 보자.

\begin{dhkcode}
print("A")

class Person:
    print("B")

print("C")
\end{dhkcode}

실행 결과:

\begin{dhkzsh}
A
B
C
\end{dhkzsh}

왜 이런 결과가 발생할까?

\zsh{class Person:}아래의 본문은 인스턴스를 생성할 때 실행되는 것이 아니라, 클래스 정의문 자체가 실행될 때 실행되기 때문이다.

실행 흐름을 살펴보자.

\begin{dhksmry}[centered]
Python 프로그램 실행
\da
print("A")
\da
class Person: 진입
\da
클래스 본문 실행
\da
print("B")
\da
클래스 객체 생성
\da
Person 이름에 클래스 객체 바인딩
\da
print("C")
\end{dhksmry}

따라서 클래스 본문에 작성된 일반적인 문장은 클래스 정의 시점에 실행된다.

반면 다음 코드는 조금 다르다.

\begin{dhkcode}
class Person:

    print("클래스 본문 실행")

    def introduce(self):
        print("메서드 실행")


p = Person()
\end{dhkcode}

실행 결과:

\begin{dhkzsh}
클래스 본문 실행
\end{dhkzsh}

\zsh{Person()}을 호출했지만 \zsh{introduce()} 내부의 출력문은 실행되지 않았다.

왜냐하면 \zsh{def introduce(self):}는 함수 본문을 즉시 실행하는 것이 아니라, 함수 객체를 만들어 클래스의 네임스페이스에 등록하는 과정이기 때문이다.

\zsh{introduce()}의 본문은 해당 메서드가 호출될 때 실행된다.

\begin{dhkcode}
p.introduce()
\end{dhkcode}

실행 결과:

\begin{dhkzsh}
메서드 실행
\end{dhkzsh}

여기서 함수 정의와 클래스 정의 사이의 차이점이 드러난다.

\begin{dhktable}
|구분|	정의문이 실행될 때|	함수·메서드 본문 실행 시점|
|:---|:---|:---|
|def|	함수 객체 생성 및 이름 바인딩|	함수 호출 시|
|class|	클래스 본문 실행 후 클래스 객체 생성 및 이름 바인딩|	클래스 본문의 메서드는 별도 호출 시|
\end{dhktable}
클래스 본문 실행과 인스턴스 생성은 서로 다른 사건이라는 점을 구분해야 한다.

\subsubsection{클래스 객체는 어떻게 만들어지는가?}

이번에는 클래스 객체가 만들어지는 과정을 조금 더 자세히 살펴보자.

다음 코드를 사용하겠다.

\begin{dhkcode}
class Person:

    species = "Human"

    def introduce(self):
        print("Hello")
\end{dhkcode}

Python은 이 클래스 정의를 처리할 때 대략 다음 과정을 거친다.

\begin{dhknote}
\begin{dhksmry}[centered]
클래스 본문에서 실행하기 위한 네임스페이스 준비
\end{dhksmry}
\end{dhknote}

클래스는 본문에서 정의되는 이름들을 관리할 공간을 준비한다.

\begin{dhknote}
\begin{dhksmry}[centered]
클래스 본문 실행
\end{dhksmry}
\end{dhknote}

\begin{dhkcode}
species = "Human"
\end{dhkcode}

문자열 객체 \zsh{"Human"}을 \zsh{species}라는 이름에 연결한다.

\begin{dhkcode}
def introduce(self):
    print("Hello")
\end{dhkcode}

함수 객체를 생성하고 \zsh{introduce}라는 이름에 연결한다.

\begin{dhknote}
\begin{dhksmry}[centered]
클래스 객체 생성
\end{dhksmry}
\end{dhknote}

클래스 본문의 실행 결과와 클래스 이름, 상속 정보 등을 사용해 클래스 객체를 생성한다.

\begin{dhknote}
\begin{dhksmry}[centered]
클래스 객체를 Person이라는 이름에 바인딩
\end{dhksmry}
\end{dhknote}

생성된 클래스 객체가 현재 네임스페이스의 \zsh{Person}이라는 이름에 연결된다.

이 과정을 개념적으로 나타내면 다음과 같다.

\begin{dhksmry}
                class Person:
                    species = "Human"

                    def introduce(self):
                        print("Hello")

                            \da

                    클래스 본문 실행

                            \da
                    
                    클래스 네임 스페이스
                        species   \ra "Human"
                        introduce \ra 함수 객체

                            \da
                    
                    클래스 객체 생성

                            \da

                    Person \ra 클래스 객체
\end{dhksmry}

실제 클래스 생성 과정에는 메타클래스, 상속 처리, 클래스 네임스페이스 준비 등의 절차가 더 포함되지만, 지금은 클래스 본문 실행 → 클래스 객체 생성 → 이름 바인딩이라는 구조를 이해하는 것이 중요하다.

\begin{dhknote}
\begin{dhksmry}[centered]
클래스 네임스페이스 확인하기
\end{dhksmry}
\end{dhknote}

클래스 객체가 생성된 뒤에는 다음과 같이 클래스의 속성을 확인할 수 있다.

\begin{dhkcode}
class Person:

    species = "Human"

    def introduce(self):
        print("Hello")


print(Person.__dict__)
\end{dhkcode}

\zsh{Person.__dict__}에는 클래스가 관리하는 속성들이 들어 있다.

개념적으로는 다음과 같은 구조다.

\begin{dhkzsh}
Person.__dict__

{
    "__module__": "__main__",
    "species": "Human",
    "introduce": <function ...>,
    "__dict__": ...,
    "__weakref__": ...,
    "__doc__": None
}
\end{dhkzsh}

실제 출력 형태는 Python 버전과 실행 환경에 따라 달라질 수 있으며, 일반적인 클래스의 \zsh{__dict__}는 읽기 전용 매핑 뷰인 \zsh{mappingproxy}로 제공된다.

여기서 중요한 것은 \zsh{species}와 \zsh{introduce}가 클래스의 속성으로 존재한다는 사실이다.

특히 \zsh{introduce}는 함수 객체라는 점을 기억하자.

\begin{dhkcode}
print(Person.__dict__["introduce"])
\end{dhkcode}

이 표현식은 클래스 네임스페이스에 저장된 함수 객체를 확인한다.

이후 2.5에서는 이 함수 객체가 인스턴스를 통해 접근될 때 어떻게 바운드 메서드로 연결되는지 살펴볼 것이다.

\subsubsection{클래스 객체도 이름을 바꾸거나 다른 변수에 전달할 수 있다}

이전에 함수도 일급 객체라는 개념을 공부했다.

\begin{dhkcode}
def add(a, b):
    return a + b

another_add = add
\end{dhkcode}

이 코드는 함수 객체를 복사하는 것이 아니라, 동일한 함수 객체를 다른 이름으로도 참조하도록 만드는 것이다.

클래스 객체도 마찬가지다.

\begin{dhkcode}
class Person:
    pass

Human = Person

print(Human is Person)
\end{dhkcode}

실행 결과:

\begin{dhkzsh}
True
\end{dhkzsh}

여기서 \zsh{Human}과 \zsh{Person}은 동일한 클래스 객체를 참조한다.

\begin{dhkdiagram}[in=6]
node p (1,1) :: Person
node c (2,2) :: 클래스 객체
node h (3,1) :: Human

edge p -> c :: HV
edge h -> c :: HV
\end{dhkdiagram}

따라서 다음 코드도 가능하다.

\begin{dhkcode}
p1 = Person()
p2 = Human()

print(type(p1) is type(p2))
\end{dhkcode}

실행 결과:

\begin{dhkzsh}
True
\end{dhkzsh}

두 이름은 같은 클래스 객체를 참조하므로 같은 클래스의 인스턴스를 생성한다.

\begin{dhknote}
\begin{dhksmry}[centered]
클래스 객체를 함수의 인수로 전달하기
\end{dhksmry}
\end{dhknote}

클래스 객체는 함수의 인수로도 전달할 수 있다.

\begin{dhkcode}
class Person:
    pass

def create_instance(cls):
    return cls()

p = create_instance(Person)

print(type(p))
\end{dhkcode}

실행 결과:

\begin{dhkzsh}
<class '__main__.Person'>
\end{dhkzsh}

여기서 실행 흐름을 분석해 보자.

\begin{dhknote}
\begin{dhksmry}[centered]
Person 평가
\end{dhksmry}
\end{dhknote}

\begin{dhkcode}
create_instance(Person)
\end{dhkcode}

\zsh{Person}이라는 이름을 평가해 클래스 객체를 가져온다.

\begin{dhknote}
\begin{dhksmry}[centered]
매개변수 cls에 바인딩
\end{dhksmry}
\end{dhknote}

\begin{dhkcode}
def create_instance(cls):
\end{dhkcode}

\zsh{cls}는 전달된 클래스 객체를 참조한다.

\begin{dhknote}
\begin{dhksmry}[centered]
클래스 객체 호출
\end{dhksmry}
\end{dhknote}

\begin{dhkcode}
return cls()
\end{dhkcode}

\zsh{cls}가 참조하는 클래스 객체를 호출해 인스턴스를 생성한다.

\begin{dhknote}
\begin{dhksmry}[centered]
인스턴스 반환
\end{dhksmry}
\end{dhknote}

생성된 인스턴스가 반환되어 \zsh{p}에 바인딩된다.

이 구조는 이후 머신러닝 프로젝트에서 유용하게 활용된다.

예를 들어 모델 클래스를 함수에 전달해 원하는 모델을 생성할 수 있다.

\begin{dhkcode}
def build_model(model_class, *args, **kwargs):
    return model_class(*args, **kwargs)
\end{dhkcode}


이 함수는 특정 모델 클래스에 종속되지 않는다.

전달받은 클래스 객체를 호출해 모델 인스턴스를 생성한다.

즉, 클래스 객체 자체를 데이터처럼 전달하고 사용할 수 있다.

\subsubsection{\zsh{type}은 무엇인가?}

앞에서 다음 결과를 확인했다.

\begin{dhkcode}
class Person:
    pass

print(type(Person))
\end{dhkcode}

실행 결과:

\begin{dhkzsh}
<class 'type'>
\end{dhkzsh}

여기서 자연스럽게 다음 질문이 발생한다.

클래스 객체의 타입이 왜 \zsh{type}일까?

Python에서 \zsh{type}은 일반적인 클래스 객체를 생성하는 기본 메타클래스(Metaclass)다.

메타클래스는 클래스 객체의 생성과 동작을 정의하는 클래스라고 이해하면 된다.

다음 관계를 살펴보자.

\begin{dhkcode}
class Person:
    pass

p = Person()

print(type(p))
print(type(Person))
print(type(type))
\end{dhkcode}

\begin{dhkzsh}
<class '__main__.Person'>
<class 'type'>
<class 'type'>
\end{dhkzsh}

이 관계를 구조적으로 표현하면 다음과 같다.

\begin{dhkdiagram}[in=12]
node t (1,2) :: type
node p (2,2) :: Person
node pp (3,2) :: p

edge t->p :: 인스턴스 관계
edge p ->pp :: 인스턴스 관계
\end{dhkdiagram}

여기서 화살표는 상속 관계가 아니라 타입과 인스턴스의 관계를 나타낸다.

\begin{dhksmry}
    - p는 Person의 인스턴스다.

    - Person은 type의 인스턴스다.

    - type 자체도 type의 인스턴스다.
\end{dhksmry}

마지막 관계는 처음 보면 조금 특이할 수 있다.

하지만 Python의 객체 모델에서는 \zsh{type}이 자기 자신을 타입으로 갖는 구조가 성립한다.

\begin{dhknote}
\begin{dhksmry}[centered]
type으로 클래스 객체를 직접 생성하기
\end{dhksmry}
\end{dhknote}

\zsh{type}은 인수를 하나 전달하면 객체의 타입을 확인하는 용도로 사용할 수 있다.

\begin{dhkcode}
type(10)
\end{dhkcode}

하지만 세 개의 인수를 전달하면 클래스 객체를 생성하는 데 사용할 수도 있다.

\begin{dhkcode}
Person = type(
    "Person",
    (),
    {"species": "Human"}
)
\end{dhkcode}

이 코드는 기본적인 경우 다음 클래스 정의와 유사한 결과를 만든다.

\begin{dhkcode}
class Person:
    species = "Human"
\end{dhkcode}

세 인수의 의미는 다음과 같다.

\begin{dhktable}
|인수|	의미|
|:---|:---|
|"Person"|	생성할 클래스의 이름|
|()|	상속할 기본 클래스들의 튜플|
|{"species": "Human"}|	클래스의 속성을 담은 딕셔너리|
\end{dhktable}

실제로 확인해 보자.

\begin{dhkcode}
Person = type(
    "Person",
    (),
    {"species": "Human"}
)

print(Person)
print(type(Person))
print(Person.species)
\end{dhkcode}

실행 결과:

\begin{dhkzsh}
<class '__main__.Person'>
<class 'type'>
Human
\end{dhkzsh}

이것은 \zsh{class}문이 단순히 텍스트 수준의 선언이 아니라, 실행 결과로 클래스 객체를 만들어 낸다는 사실을 보여주는 예제다.

다만 일반적인 \zsh{class}문과 \zsh{type()}을 직접 호출하는 방식이 모든 상황에서 완전히 동일한 것은 아니다.

클래스 정의문에는 클래스 본문 실행, 메타클래스 선택, \zsh{__class__} 참조 처리 등의 추가적인 의미가 있기 때문이다.

현재 단계에서는 메타클래스를 직접 작성하는 방법까지 공부할 필요는 없다.

우선 다음 관계만 이해하면 충분하다.

\begin{dhkdiagram}[in=6]
node i (1,1) :: Instance
node c (1,3) :: Class
node m (1,5) :: Metaclass

edge i -> c :: type
edge c -> m :: type
\end{dhkdiagram}

\subsubsection{함수 객체와 클래스 객체의 공통점과 차이}

이번에는 함수 객체와 클래스 객체를 비교해 보자.

\begin{dhknote}
\begin{dhksmry}[centered]
함수 정의
\end{dhksmry}
\end{dhknote}

\begin{dhkcode}
def add(a, b):
    return a + b
\end{dhkcode}

이 코드를 실행하면 함수 객체가 생성되고 \zsh{add}라는 이름에 바인딩된다.

\begin{dhksmry}
    add \ra 함수 객체
\end{dhksmry}

\begin{dhknote}
\begin{dhksmry}[centered]
클래스 정의
\end{dhksmry}
\end{dhknote}

\begin{dhkcode}
class Calculator:
    pass
\end{dhkcode}

이 코드를 실행하면 클래스 객체가 생성되고 \zsh{Calculator}라는 이름에 바인딩 된다.

\begin{dhksmry}
    Calculator \ra 클래스 객체
\end{dhksmry}

두 객체 모두 Python 프로그램 안에서 참조하거나 다른 함수에 전달할 수 있다.

하지만 호출했을 때의 기본적인 역할은 다르다.

\begin{dhkcode}
add(1, 2)
\end{dhkcode}

함수 객체를 호출하면 함수 본문이 실행된다.

반면 다음 코드에서는

\begin{dhkcode}
Calculator()
\end{dhkcode}

클래스 객체를 호출해 일반적으로 새로운 인스턴스를 생성하고 초기화한다.

정리하면 다음과 같다.

\begin{dhktable}
|구분|	함수 객체|	클래스 객체|
|:---|:---|:---|
|정의 문법|	def	|class|
|정의 결과|	함수 객체	|클래스 객체|
|일반적인 타입|	function	|type|
|이름에 바인딩|	가능	|가능|
|인수로 전달|	가능	|가능|
|반환값으로 사용|	가능	|가능|
|() 호출|	함수 본문 실행	|일반적으로 인스턴스 생성 및 초기화|
|주요 역할|	연산과 동작 수행	|인스턴스의 구조와 //*동작 정의|
\end{dhktable}

여기서 주의할 점이 있다.

함수와 클래스 모두 호출 가능한 객체라는 공통점이 있지만,

모든 Callable 객체가 함수이거나 클래스인 것은 아니다.

예를 들어 \zsh{__call__()}을 구현한 인스턴스 역시 Callable 객체가 될 수 있다.

이 내용은 2.10에서 다시 연결할 것이다.

\subsubsection{PyTorch와의 연결 - 모델 클래스도 객체다}

이제 처음에 살펴봤던 PyTorch 코드로 돌아가 보자.

\begin{dhkcode}
import torch.nn as nn


class MyModel(nn.Module):

    def __init__(self):
        super().__init__()
        self.linear = nn.Linear(10, 5)

    def forward(self, x):
        return self.linear(x)
\end{dhkcode}

여기서 다음 표현식에 주목하자.

\begin{dhkcode}
MyModel
\end{dhkcode}

이것은 모델 인스턴스가 아니다.

모델 클래스를 나타내는 클래스 객체다.

반면 다음 코드는

\begin{dhkcode}
model = MyModel()
\end{dhkcode}

\zsh{MyMdodel} 클래스 객체를 호출해 모델 인스턴스를 생성한다.

따라서 다음과 같이 구분할 수 있다.

\begin{dhkbranch}[in=8]
node main :: MyModel
arrow main
node main :: 모델 클래스 객체
arrow main :: MyModel() 호출
node main :: 모델 인스턴스 객체
arrowup main
node main :: model
\end{dhkbranch}

이 차이는 머신러닝 코드를 작성할 때 중요하다.

예를 들어 다음 함수를 보자.

\begin{dhkcode}
def create_model(model_class):
    return model_class()
\end{dhkcode}

이 함수에는 클래스 객체를 전달한다.

\begin{dhkcode}
model = create_model(MyModel)
\end{dhkcode}

반대로 이미 생성된 모델 인스턴스를 전달하는 경우도 있다.

\begin{dhkcode}
def count_parameters(model):
    return sum(
        p.numel()
        for p in model.parameters()
    )
\end{dhkcode}

이 함수는 모델 클래스가 아니라 모델 인스턴스를 전달받는 것을 전제로 한다.

\begin{dhkcode}
model = MyModel()

count = count_parameters(model)
\end{dhkcode}

따라서 PyTorch 코드를 읽을 때는 다음을 구분해야 한다.

\begin{dhksmry}
    1. 클래스 객체를 전달하는가?

    2. 클래스 객체를 호출해 인스턴스를 생성하는가?

    3. 이미 생성된 인스턴스를 전달하는가?

    4. 인스턴스를 Callable 객체로 호출하는가?
\end{dhksmry}

이 구분이 가능해지면 Hugging Face나 PyTorch에서 클래스를 인수로 전달하거나 모델을 동적으로 생성하는 코드도 훨씬 이해하기 쉬워질 것이다.

\subsubsection{핵심 정리}

이번에는 Python에서 클래스도 객체라는 사실을 중심으로 \zsh{class} 정의문의 실행 과정을 살펴보았다.

핵심 내용을 정리하면 다음과 같다.

\begin{dhknote}
\begin{dhksmry}[centered]
클래스도 Python 객체다.
\end{dhksmry}
\end{dhknote}

\begin{dhkcode}
class Person:
    pass
\end{dhkcode}

이 코드를 실행하면 클래스 객체가 생성되고 \zsh{Person}이라는 이름에 바인딩된다.

\begin{dhknote}
\begin{dhksmry}[centered]
클래스 정의문은 실행된다.
\end{dhksmry}
\end{dhknote}

클래스 본문은 클래스 정의 시점에 실행되며, 그 결과가 클래스 객체를 구성하는 데 사용된다.

\begin{dhknote}
\begin{dhksmry}[centered]
클래스 객체와 인스턴스 객체는 다르다.
\end{dhksmry}
\end{dhknote}

\begin{dhkcode}
Person      # 클래스 객체
Person()    # 일반적으로 인스턴스 생성
\end{dhkcode}

\begin{dhknote}
\begin{dhksmry}[centered]
클래스 객체도 다른 이름에 바인딩하거나 함수에 전달할 수 있다.
\end{dhksmry}
\end{dhknote}

\begin{dhkcode}
Huamn = Person
\end{dhkcode}

이 경우 두 이름은 동일한 클래스 객체를 참조한다.

\begin{dhknote}
\begin{dhksmry}[centered]
일반적인 클래스 객체의 타입은 type이다.
\end{dhksmry}
\end{dhknote}

\begin{dhkcode}
type(Person)
\end{dhkcode}

\zsh{type}은 Python에서 기본적인 메타클래스 역할을 한다.

\begin{dhknote}
\begin{dhksmry}[centered]
클래스 객체를 호출하면 일반적으로 인스턴스 생성 과정이 시작된다.
\end{dhksmry}
\end{dhknote}

\begin{dhkcode}
p = Person()
\end{dhkcode}

인스턴스 생성과 초기화의 구체적인 과정은 이후 \zsh{__new__()}와 \zsh{__init__()}을 학습하면서 자세히 살펴본다.

\begin{dhknote}
\begin{dhksmry}[centered]
PyTorch의 모델 클래스 역시 클래스 객체다.
\end{dhksmry}
\end{dhknote}

\begin{dhkcode}
class MyModel(nn.Module):
    ...
\end{dhkcode}

여기서 \zsh{MyModel}은 클래스 객체이며, \zsh{MyModel()}을 호출하면 일반적으로 모델 인스턴스가 생성된다.

\begin{dhknote}
\begin{dhksmry}[centered]
반드시 이해해야 할 결론
\end{dhksmry}
\end{dhknote}

이전에 함수 객체를 공부하며 다음 구조를 이해했다.

\begin{dhksmry}[centered]
def \ra Function Object \ra Name Binding
\end{dhksmry}

이번에는 여기에 대응하는 새로운 구조를 배웠다.

\begin{dhksmry}[centered]
class \ra Class Body Execution \ra Class Object \ra Name Binding
\end{dhksmry}

즉, Python의 \zsh{class}는 클래스 객체를 생성하는 실행 가능한 문장이며,

생성된 클래스 객체는 다른 Python 객체들과 마찬가지로 이름에 바인딩하고 전달하고 호출할 수가 있다.

이것이 이번 학습에서 가장 중요한 개념이다.

다음 \zsh{2.3. 클래스와 인스턴스}에서는

이 개념을 기반으로 \zsh{Person}과 \zsh{Person()}이 구체적으로 어떻게 다른지, 

클래스 객체를 호출했을 때 인스턴스가 생성되고 이름에 바인딩되는 과정을 더 자세히 살펴볼 것이다.
